\documentclass[11pt,a4paper]{article}
\usepackage[margin=25mm]{geometry}
\usepackage{amsmath,amssymb,amsthm,booktabs,longtable,hyperref}
\setlength{\emergencystretch}{2em}
\newtheorem{theorem}{Theorem}\newtheorem{lemma}[theorem]{Lemma}
\newcommand{\PSL}{\operatorname{PSL}}\newcommand{\PGL}{\operatorname{PGL}}
\newcommand{\Aut}{\operatorname{Aut}}\newcommand{\F}{\mathbb F}
\title{Local actions in rank two over a projective line}
\author{Human Compression project}\date{}
\begin{document}\maketitle
We reprove the principal classification in Part I of De Saedeleer and
Leemans, \emph{On the rank two geometries of the groups $\PSL(2,q)$},
Ars Mathematica Contemporanea 3 (2010), 177--192. Throughout, $q\ge4$
is a prime power, $d=\gcd(2,q-1)$ and $G=\PSL(2,q)$. The small
nonsimple groups are not included in this standing hypothesis.
We retain Dickson's classification of the proper subgroups of $G$, including
the usual maximal-subgroup and natural-subfield clauses. The projective-line
action, torus normalizers and the orders of $\PSL$ and $\PGL$ are classical
inputs. We do not retain the source's principal geometry tables as premises.

The reduction below uses two elementary facts about primitive actions.
It leaves two uniform families and fourteen small fields. Exact subgroup
intersections and graph computations then determine every entry of the
finite list. We distinguish ordered parabolics from the unordered convention
stated on page 181 of the source. Two conjugacy counts differ between those
conventions. The corrections to counts and diagram labels are recorded
explicitly at the end.

\section{The local action test}
For distinct proper subgroups $H,K<G$, let $B=H\cap K$. Vertices of the
two types are the right cosets of $H$ and $K$; incidence means nonempty
intersection. The neighbours of the base $H$-vertex are indexed by $H/B$,
and similarly at the $K$-vertex. Every incident pair is a translate of the
base pair. Thus the geometry is flag-transitive, and is connected precisely
when $\langle H,K\rangle=G$.

\begin{lemma}\label{local}
The rank-two system is locally doubly transitive and residually weakly
primitive if and only if at least one of $H,K$ is maximal in $G$ and
\[
 |B\backslash H/B|=|B\backslash K/B|=2.
\]
These conditions imply firmness and connectedness. For fixed $H$, conjugacy
classes of ordered systems are the $N_G(H)$-orbits on the eligible $K$;
the analogous $\Aut(G)$-classes use $\Aut(G)_H$. For unordered systems,
one must additionally identify the orbits transported by interchanging $H,K$.
\end{lemma}
\begin{proof}
The orbits of a point stabilizer $B$ on $H/B$ are precisely its double
cosets. Exactly two orbits means that the nonbase points form one orbit,
so the action is doubly transitive and its degree is at least two. It is
primitive, whence $B$ is maximal in each parabolic. Conversely these local
actions give every nonempty-residue condition; the empty-residue condition
is that one parabolic be maximal in $G$. The other parabolic is not contained
in it, since their intersection has proper index in both. They therefore
generate $G$. The orbit statements follow by normalizing the first subgroup.
If a transporter sends $K$ to $H$, its image of $H$ gives the additional
identification. Two choices of transporter differ by the indicated stabilizer.
\end{proof}

\begin{lemma}\label{primitive}
In a faithful primitive action, a nontrivial abelian normal subgroup is
regular. A doubly transitive action with cyclic regular normal subgroup has
prime degree; if its point stabilizer has order at most two, that degree
is two or three. In every doubly transitive action of $L$ on $L/B$,
\[
 [L:B]-1\ \mid\ |B|.
\]
\end{lemma}
\begin{proof}
The orbits of a normal subgroup form blocks, so a nontrivial normal subgroup
in a faithful primitive action is transitive. If it is abelian, all its point
stabilizers coincide and lie in the kernel, proving regularity. For a cyclic
regular normal subgroup, any proper nontrivial characteristic subgroup would
again be normal and transitive, which is impossible. Its order is therefore
prime. The point stabilizer is transitive on all other points, proving both
the bound on the degree and the divisibility assertion.
\end{proof}

\section{The two uniform families}
Write $P=V:C_t$ for a Borel subgroup, where $V=(\F_q,+)$ and
$t=(q-1)/d$. A generator of its scalar complement generates the field
$\F_q$: indeed $t>\sqrt q-1$, whereas any proper subfield has multiplicative
order at most $\sqrt q-1$. Consequently its only invariant additive
subgroups are $0,V$.

\begin{theorem}\label{borel}
If one parabolic is a Borel subgroup, then $q$ is even. There are exactly
two systems up to conjugacy, hence also up to $\Aut(G)$:
\begin{enumerate}
\item two copies of the projective line, with unequal points incident;
\item the points and unordered pairs of points of the projective line,
with incidence given by containment.
\end{enumerate}
Their local arc-transitivity depths are respectively two and three.
Their type sizes, local degrees, gonalities and type diameters are
\[
\begin{array}{c|c|c|c|c}
 &\text{type sizes}&\text{local degrees}&\text{gonality}&\text{diameters}\\\hline
1&(q+1,q+1)&(q,q)&2&(3,3)\\
2&(q+1,q(q+1)/2)&(q,2)&3&(3,4).
\end{array}
\]
The second is a partial linear space; the first is not. Both are residually
primitive. The first is the usual bipartite double of the complete graph
underlying the second.
\end{theorem}
\begin{proof}
In a primitive coset action of $P$, the kernel intersected with $V$ is
scalar-invariant. If the image of $V$ is trivial, the action factors through
the cyclic complement and double transitivity forces degree two. Thus the
point stabilizer contains $V$. No proper subgroup of $G$ outside this Borel
contains $V$: the subgroup alternatives show that a proper subfield group
has a strictly smaller characteristic Sylow subgroup, while the small
exceptional groups give no exception for $q\ge4$. In the remaining affine
case the unique fixed projective point forces the same Borel. Hence a second
parabolic is impossible in this case.

Otherwise $V$ acts faithfully and regularly, by Lemma~\ref{primitive},
and the degree is $q$. Double transitivity requires the scalar complement
to be transitive on the nonzero translations. It follows that $t=q-1$,
so $q$ is even. The point stabilizer is a complement $C_{q-1}$: a generator
of its projection is an affine map with a unique fixed point, so translation
conjugacy puts the whole complement in scalar form.

A cyclic torus $C_{q-1}$ fixes two projective points. A proper overgroup
acting doubly transitively on its cosets is one of their two Borels or its
dihedral normalizer. To check this directly from Dickson, an affine overgroup
has regular translation group and a scalar complement of order one less
than its size, forcing the full Borel. A cyclic subgroup cannot be maximal
in a genuine subfield $\PSL$ or $\PGL$ group except in the small $S_3,A_4$
cases. The exceptional groups likewise supply only $A_4$ with stabilizer
$C_3$, which is the Borel at $q=4$. A dihedral overgroup normalizes the
torus and is therefore its full normalizer. This proves completeness.

The two stabilizer configurations give exactly the stated incidence models.
The projective action is doubly transitive, so they each form one conjugacy
class. For even $q$, $G=\PGL(2,q)$ acts sharply transitively on ordered
triples of distinct projective points: the fractional linear transformation
carrying a given triple to $(\infty,0,1)$ is unique. In the first graph,
two-arcs starting at a fixed point correspond to distinct triples. At length
three the last point may equal the first underlying point or not, so these
arcs cannot form one orbit. In the point-pair incidence graph, arcs alternate
between points and edges. Sharp transitivity on distinct triples gives
local transitivity through length three. Starting at an edge, length-four
arcs distinguish return to its other endpoint from a new endpoint, so the
depth is exactly three. Direct paths in the two incidence models give the
displayed gonalities and diameters. Two vertices in the first part of model~1
have $q-1\ge3$ common neighbours, excluding the partial-linear property;
it is not a digon either. Model~2 is the incidence geometry of a complete
graph. The classical maximality of these Borels and split torus normalizers
proves residual primitivity.
\end{proof}

\section{Why only small exceptional fields remain}
Let $H=D_{2n}$ be a full torus normalizer, where $n=(q\pm1)/d$.

\begin{lemma}\label{dihedral}
If $H$ acts doubly transitively on $H/B$, then either $B=C_n$, or
$B=D_{2t}$ with $n=2t$ or $n=3t$. Apart from the uniform families in
Theorem~\ref{borel}, every eligible system occurs for
\[
q\in\{4,5,7,8,9,11,13,16,17,19,23,25,29,31\}.
\]
\end{lemma}
\begin{proof}
Pass to the faithful quotient of the primitive action. The cyclic rotation
subgroup either has trivial image, giving the two-point quotient with
$B=C_n$, or is regular of prime order. The image of the point stabilizer
then has order at most two. Lemma~\ref{primitive} gives degrees two or
three, and the remaining subgroups are exactly the indicated dihedral ones.

First let $B=C_n$, with $n>2$. Any dihedral overgroup normalizes this
semisimple cyclic subgroup and lies in $N_G(B)=H$, so cannot be the second
parabolic. Cyclic overgroups likewise give none. Suppose the second parabolic
is affine, with normal elementary abelian characteristic subgroup of order
$p^e$. Since $p\nmid n$, its image is regular with trivial kernel. The cyclic
scalar stabilizer acts freely on its nonzero vectors and transitively on
them, so $n=p^e-1$, with $e\le r$ when $q=p^r$. For odd $p$ the split
equation $2p^e=p^r+1$ is impossible modulo $p$. The nonsplit equation is
$2p^e=p^r+3$; it forces $p=3$ and then $e=r=1$, outside $q\ge4$.
For even $p$ only the split solution $e=r$ remains, giving the full Borel.
In the exceptional groups a maximal cyclic point stabilizer can only be
$C_3<A_4$. In a subfield $\PSL$ or $\PGL$ group, a cyclic subgroup lies
in a proper torus normalizer unless it belongs to a Borel; that Borel is
cyclic only in the small $S_3,A_4$ cases. These are already covered, or give
$n=3$. The cases $n\le3$ have $q\le7$ and are included in the finite list.

Now suppose $B=D_{2t}$ and $t>2$. Its characteristic rotation subgroup
$C_t$ is semisimple. A second dihedral parabolic would normalize $C_t$,
and again lie in $H$, which is impossible. If a second affine parabolic
has normal characteristic subgroup $V$, then $V\cap B=1$. In odd
characteristic this follows from orders. In characteristic two, $t$ is
odd and $D_{2t}$ has no nontrivial normal $2$-subgroup. The regular normal
image of $V$ would therefore identify $B$ with a cyclic complement,
contrary to $t>2$.

Consider a second subfield group $L=\PSL(2,s)$ or $\PGL(2,s)$.
Because $B$ is maximal in $L$, the classical subgroup list makes it a full
torus normalizer, unless $L$ is one of the small $S_3,A_4,S_4,A_5$ cases.
Indeed a dihedral Borel is possible only at these small fields, and a
dihedral subfield subgroup itself reduces to $S_3$. In odd $\PSL$,
$|B|\le s+1$ and $[L:B]\ge s(s-1)/2$, so Lemma~\ref{primitive}
forces $s\le4$. In even $\PSL$ or in $\PGL$, the corresponding bounds
are $|B|\le2(s+1)$ and $[L:B]\ge s(s-1)/2$, forcing $s\le6$.
The only additional possibility is $\PGL(2,5)$: its torus normalizers
have orders $8,12$ and indices $15,10$, respectively, and neither
$14\mid8$ nor $9\mid12$. It supplies no case.

The dihedral doubly transitive point stabilizers in the remaining exceptional
groups are $S_3,D_8<S_4$ and $D_{10}<A_5$. This follows from their
elementary subgroup lists and the same divisibility test; $A_4$ supplies
only its cyclic $C_3$. Thus $t\in\{3,4,5\}$ and $n\in\{2t,3t\}$.
Solving $q\pm1=dn$ among prime powers gives the displayed finite fields.
Finally $t\le2$ gives $n\le6$ and adds only $q\le13$ from that list.
We include the nonmaximal torus cases $q=5,7,9,11$ explicitly, rather than
assuming that the specified torus parabolic is the maximal one.
\end{proof}

\section{The finite list and its verification}
Table~\ref{finite} removes repetitions of the two uniform families and the
isomorphic ambient groups $\PSL(2,4)\cong\PSL(2,5)$. Here $c$ counts
ordered conjugacy classes; $u$ counts unordered conjugacy classes. Every row
is a single $\Aut(G)$-class. Type sizes and degrees follow from
$|G|/|H|,|G|/|K|$ and $|H|/|B|,|K|/|B|$. The last columns give
gonality, the two type diameters, and local arc-transitivity depth.
\begin{table}[ht]
\centering\small
\begin{tabular}{rlllrrlr}
\toprule
$q$&$H$&$K$&$B$&$c$&$u$&$(g;d_0,d_1)$&$s$\\\midrule
$4$&$S_3$&$C_2^2$&$C_2$&1&1&$(5;5,6)$&3\\
$4$&$S_3$&$S_3$&$C_2$&1&1&$(3;5,5)$&2\\
$11$&$D_{12}$&$D_{12}$&$C_2^2$&2&1&$(5;7,7)$&3\\
$13$&$D_{12}$&$D_{12}$&$C_2^2$&2&1&$(6;9,9)$&3\\
$16$&$D_{30}$&$A_5$&$D_{10}$&1&1&$(4;7,7)$&3\\
$17$&$D_{16}$&$S_4$&$D_8$&2&2&$(9;14,14)$&7\\
$17$&$D_{18}$&$S_4$&$S_3$&2&2&$(4;8,8)$&3\\
$19$&$D_{20}$&$A_5$&$D_{10}$&2&2&$(5;8,7)$&3\\
$23$&$D_{24}$&$S_4$&$D_8$&2&2&$(6;11,11)$&3\\
$25$&$D_{24}$&$S_4$&$D_8$&2&2&$(5;13,13)$&3\\
$29$&$D_{30}$&$A_5$&$D_{10}$&2&2&$(4;8,7)$&3\\
$31$&$D_{30}$&$A_5$&$D_{10}$&2&2&$(4;8,7)$&3\\
\bottomrule
\end{tabular}
\caption{The exceptional systems. The two $q=4$ rows also occur at $q=5$;
the $A_4,S_3,C_3$ row at $q=5$ is the $q=4$ uniform point-pair geometry.}
\label{finite}\end{table}

\begin{theorem}\label{main}
For $q\ge4$, Theorem~\ref{borel} and Table~\ref{finite} give all locally
doubly transitive RWPRI rank-two systems having a Borel or full torus
normalizer as one parabolic. The finite rows are partial linear spaces.
Both parabolics are maximal except in the $q=4$, $S_3,C_2^2$ row and
the $q=25$ row; precisely the other finite rows are residually primitive.
\end{theorem}
\begin{proof}
The two preceding results reduce completeness to the fourteen fields.
The certificate constructs $G$ in its natural permutation action and lists
all its subgroups from conjugacy-class representatives and their conjugates.
For each representative full torus normalizer $H$, it tests every proper
$K$ by literal intersection, the two double-coset counts, and actual maximality
of at least one parabolic. No abstract order match substitutes for the
intersection or action test. Lemma~\ref{local} proves that precisely the
eligible systems survive. The normalizer and automorphism-stabilizer orbits
give their ordered counts. Actual transporters implement the unordered
identifications. The full returned inventory is bound to the certificate,
including the empty cases and the uniform rows removed from the display.

For each representative, right cosets of $B$ supply exactly $|G:B|$ edges
between its two parabolic coset sets. The permutation action on the graph
is faithful. Breadth-first distances from a vertex of each type give the
diameters, since $G$ is transitive on each type. The girth is computed both
by deleting an edge and finding the distance between its endpoints, and by
an independent all-root breadth-first computation. All finite rows have
girth at least six, proving the partial-linear property.

For completeness, local arc-depth is computed as follows. Start at a vertex
and its stabilizer. If the current pointwise path stabilizer is transitive
on the neighbours of the last vertex other than its predecessor, extend by
one chosen neighbour and intersect with its stabilizer. Otherwise stop.
Inductively all arcs of the attained length form one orbit, so this test is
necessary and sufficient at the next length. Repeated vertices other than
immediate backtracking are allowed, as required by the source definition.
Perform this from both types and take the smaller depth. The exact outputs
are those in the table.
\end{proof}

\section{A scalar-subspace count}
The source's Remark~4.3 gives eleven conjugacy classes of elementary abelian
subgroups of order sixteen in $\PSL(2,64)$, in correction of an earlier
enumeration attributed there to Dickson. This count has a short independent
proof. A nontrivial unipotent subgroup fixes a unique projective point.
After fixing that point, subgroups of order sixteen are the four-dimensional
$\F_2$-subspaces $T$ of $\F_{64}$, with scalar multiplication by
$\F_{64}^{\times}$ giving precisely conjugacy. Translations centralize them.

The set $E=\{a\in\F_{64}:aT\subseteq T\}$ is a subfield: it is closed
under addition and multiplication, and a nonzero scalar acts bijectively on
the finite space $T$, so its inverse also preserves $T$. Its degree over
$\F_2$ divides both six and four. Thus $E=\F_2$ or $\F_4$.
There are $\binom{6}{4}_2=651$ such subspaces in all, and
$\binom{3}{2}_4=21$ are $\F_4$-spaces. The latter have scalar stabilizer
of order three, so form one orbit of length $63/3=21$. The remaining
$630$ subspaces have trivial scalar stabilizer and form ten orbits of
length 63. Their normalizers in $G$ have orders $64\cdot3$ and $64$,
giving one subgroup conjugacy class of length 1365 and ten of length 4095,
exactly as stated in the source. This argument counts subspaces and their
scalar orbits separately; treating the 651 subspaces as conjugacy classes
would conflate the two sets.

\section{Scope and corrections}
The first four sections reprove the main existence and completeness result,
the local arc-depth column, and the diagram invariants of the source's
Theorem~1.1 and Propositions~5.3, 5.8 and 5.12. The omitted intermediate
overgroup proofs are supplied by the two uniform action arguments and
Lemma~\ref{dihedral}. Dickson's classification and its maximal-subgroup
clauses remain attributed classical inputs. The names Petersen, Desargues
and Perkel refer to earlier identifications; no priority claim about these
graphs is made here.

The source defines conjugacy using the unordered set of parabolics. At
$q=11,13$, its two ordered $D_{12},D_{12}$ classes are exchanged by
type interchange. Thus the printed count two is correct for ordered
parabolics, and the count under its stated unordered convention is one.
The certificate records an actual transporter and checks the distinct
ordered orbits. The $\Aut(G)$ count remains one.

The split torus normalizer has index $q(q+1)/2$, as printed correctly on
page 186; the repeated $q(q-1)/2$ on page 188 is incorrect. In the
odd-characteristic Table~3, the nonsplit normalizer index is
$q(q-1)/2$, with a factor two missing from the printed denominator.
In the
$q=17$ diagram on page 189, the $D_{16}$ type has $2448/16=153$
vertices, not 136. All other displayed exceptional distances, gonalities
and local arc-depths agree with the exact computations. The unqualified
prime-order reflection clause in Table~4 is replaced by
Lemma~\ref{dihedral}: $D_{2p}$ on reflection cosets is doubly transitive
only for $p=3$ (with the degree-two boundary treated separately).

The simple-group hypothesis matters at the boundary. In $\PSL(2,2)=S_3$,
two distinct order-two parabolics give the six-cycle. The two uniform rows
coincide there, and this cycle is locally arc-transitive at every length:
after its first edge a nonbacktracking walk has a unique continuation, and
the vertex stabilizer exchanges the two choices of direction. In
$\PSL(2,3)=A_4$ no eligible system exists. A cyclic order-three parabolic
has no doubly transitive nontrivial coset action. A Klein-four parabolic
would require an order-two intersection, but every proper overgroup of that
intersection lies in the same unique Klein four. This also settles the
boundary if the displayed source theorem is read without its surrounding
simple-group convention.

Classical subgroup classifications remain borrowed inputs; the
finite certificate supplements the written reduction and is not a proof
of those classifications.
\end{document}
